% 教学反思 · 第 1 次课
\section{第 1 次课：集合的含义与表示}

\begin{teachreflect}{达成情况}
\begin{itemize}
  \item 列举法与描述法的互化，多数学生能完成；
  \item 「代表元素决定集合」这一点，约一半学生仍停留在背结论。
\end{itemize}
\end{teachreflect}

\begin{teachreflect}{意外之处}
学生把 $\{x^2-1=0\}$ 当成集合写法，且坚信竖线后可以放方程。
原以为讲完代表元素就能避免，实际需要在黑板上并排写三个集合对比才有效。
\end{teachreflect}

\begin{teachreflect}{下次调整}
\begin{enumerate}
  \item 把「三个集合对比」提前到概念之后立刻做，不留到例题；
  \item 增加一道「判断写法是否正确」的口答；
  \item 练习第 3 题（大于 3 的实数）留作课后，不占课堂时间。
\end{enumerate}
\end{teachreflect}
